\guidelinesubsubsection{Rationale}

Meaningful evaluation requires well-understood, valid measurement instruments.
Without justified benchmarks and metrics, claims about LLM performance lack the rigor needed for scientific comparison and cumulative progress.
A \textit{benchmark} is a ``standard tool for the competitive evaluation and comparison of competing systems or components according to specific characteristics, such as performance, dependability, or security''~\cite{DBLP:conf/wosp/KistowskiAHLHC15}. ``A \textit{metric} is a method, algorithm, or procedure for assigning one or more numbers to a phenomenon''~\cite{DBLP:books/sp/24/RalphKAA24}. A baseline is a reference point, enabling comparison of LLMs against traditional algorithms with lower computational costs.

\guidelinesubsubsection{Recommendations}

Benchmark and metric selection requires understanding the benchmark tasks, what exactly is being measured, and how it relates to the (often latent) variables researchers actually care about. When one or more metrics or benchmarks are used, researchers \must briefly justify in the \paper why each selected benchmark and metric is suitable for the given task or study, and \must discuss the reliability and validity (especially construct validity) of those choices. Beyond these requirements, researchers \should:
\begin{itemize}
    \item provide an operational definition of the phenomenon the benchmark is intended to measure (e.g., functional correctness, code maintainability, vulnerability detection), including its scope and any sub-components~\cite{bean2025measuring};
    \item summarize the structure and tasks of the selected benchmark(s), including the programming language(s) and descriptive statistics such as the number of contained tasks and test cases;
    \item describe and justify the sampling strategy used to select problems for inclusion in the benchmark (e.g., function-completion problems in \emph{HumanEval}, GitHub issues in \emph{SWE-bench}, vulnerable functions in \emph{DiverseVul}); if non-probability sampling (e.g., convenience) is used, discuss its implications for the generalizability of conclusions~\cite{DBLP:journals/ese/BaltesR22, bean2025measuring};
    \item discuss the limitations of the selected benchmark(s) (e.g., widely used benchmarks such as \emph{HumanEval} \cite{DBLP:journals/corr/abs-2107-03374} and \emph{MBPP} \cite{DBLP:journals/corr/abs-2108-07732} only test short Python functions, which is not representative of the full breadth of SE work~\cite{Chandra2025benchmarks});
    \item include an example of a task and the corresponding test case(s) to illustrate the structure of the benchmark.
\end{itemize}

If multiple benchmarks exist for the same task, researchers \should compare both performance and design choices (e.g., which tasks are included, how outputs are scored, what aspect of the phenomenon is covered) across benchmarks~\cite{bean2025measuring}.
When selecting only a subset of all available benchmarks, researchers \should use the most specific benchmarks given the context.
When adapting an existing benchmark, researchers \should document what was changed and why, and \should report performance on both the original and the adapted version where feasible~\cite{bean2025measuring}.

Benchmark scores often confound the target capability with unrelated capabilities the task happens to require. For code translation benchmarks, prompt formatting alone can shift performance by up to 40\%~\cite{DBLP:journals/corr/abs-2411-10541, cao2025should}, conflating prompt format sensitivity with translation capability. More broadly, output format compliance (e.g., specific JSON schemas or unit-test conventions) and general instruction-following capability are routinely bundled into the score~\cite{bean2025measuring}. Researchers \should identify which capabilities a benchmark conflates, isolate the target capability where possible (e.g., by reporting per-subtask breakdowns or by adopting benchmarks designed to test the target capability in isolation), and acknowledge remaining confounders as threats to construct validity.

After running a benchmark, researchers \should perform an error analysis by categorizing the failures observed and reporting the relative frequency of each category. If failures concentrate on inputs that demand capabilities other than the target capability (e.g., reading across many files rather than fixing the bug the benchmark targets), this is a construct-validity threat and \should be reported alongside the primary scores~\cite{bean2025measuring}.

In addition to disclosing data sources and collection dates (see \emph{Challenges} below), researchers creating or releasing a new benchmark \should adopt concrete contamination-prevention mechanisms: maintaining a held-out subset of items for ongoing, uncontaminated evaluation; embedding canary strings, that is, unique markers that downstream tools can later search for in model outputs to detect inclusion in training data; and investigating whether the benchmark's source materials may already appear in common LLM training corpora~\cite{bean2025measuring, cao2025should}. Researchers using existing benchmarks \should additionally discuss contamination as a study limitation (\limitationsmitigations).

Furthermore, researchers \should check whether a less resource-intensive approach (e.g., for static analysis tasks or program repair) can serve as a baseline. If so, the LLM or LLM-based tool \should be compared with such baselines using suitable metrics. Even if LLM-based tools outperform baselines, researchers \should discuss whether the resources consumed justify the (often marginal) improvements~\cite{DBLP:journals/cacm/Menzies25}.

To compare traditional and LLM-based approaches or different LLM-based tools, researchers \should report established metrics whenever possible, as this allows secondary research.
They can report additional metrics that they consider appropriate.
We briefly discuss common metrics in the \emph{Examples} subsection below.
As mentioned, researchers \must motivate why they chose a certain metric or variant thereof for their particular study.
Prior adoption alone does not constitute sufficient justification; researchers \mustnot justify metric choices solely by citing their use in prior work.

Due to LLM non-determinism, researchers \should repeat experiments and report descriptive statistics of model or tool performance (e.g., arithmetic mean, median, confidence intervals, standard deviations)~\cite{DBLP:conf/nips/AgarwalSCCB21, DBLP:journals/corr/abs-2602-07150}. If comparing models or tools, researchers \must use appropriate inferential statistics (e.g., hypothesis tests such as the Mann-Whitney U test, McNemar's test for binary outcomes~\cite{DBLP:conf/iclr/KublerBKZYKK26}, or bootstrap-based comparisons, along with effect sizes) rather than relying solely on differences in means or other summary statistics. When scoring uses ratings that vary across raters or runs (e.g., human raters, LLM-as-judge), researchers \should report the distribution of ratings per item rather than only aggregated point estimates or exact-match agreement, since aggregation can mask systematic disagreement~\cite{bean2025measuring}.

The number of required repetitions depends on factors such as the study type, the variability of the task, and the desired precision of estimates. As with sample sizes for human validation (\humanvalidation), researchers \should justify their chosen number of repetitions, for example, through a power analysis~\cite{Cohen1992, DBLP:journals/corr/abs-2602-07150} or by monitoring the convergence of descriptive statistics across incremental runs~\cite{DBLP:journals/corr/abs-2410-03492}. A pilot study can help estimate the expected variability and inform this decision.

From a measurement perspective, researchers \should reflect on the theories, values, and measurement models on which the benchmarks and metrics they have selected for their study are based.
For example, labeling phenomena as ``bugs'' in a large open dataset reflects a certain theory of what constitutes a bug, as well as the values and perspectives of the people who labeled the dataset. Reflecting on the context in which these labels were assigned and discussing whether and how the labels generalize to a new study context is crucial.

Researchers building or releasing new SE benchmarks \should consult operational checklists. \citeauthor{cao2025should}~\cite{cao2025should} provide \emph{HOW2BENCH}, a 55-item checklist covering benchmark design, construction, evaluation, analysis, and release. \citeauthor{bean2025measuring}~\cite{bean2025measuring} provide a complementary domain-agnostic checklist organized around the construct-validity recommendations summarized above.

\guidelinesubsubsection{Examples}

\paragraph{Common Metrics:}

Two common metrics used for generation tasks are \emph{BLEU-N} and \emph{pass@k}.
\emph{BLEU-N}~\cite{DBLP:conf/acl/PapineniRWZ02} was originally developed for evaluating machine translation quality by measuring n-gram precision between a candidate and reference text, ranging from 0 (dissimilar) to 1 (similar). It has been widely adopted in SE for code generation tasks, though its validity in this context is debatable: n-gram overlap does not capture functional correctness, and syntactically different code can be semantically equivalent (see \emph{Challenges} below).
Code-specific variations such as \emph{CodeBLEU}~\cite{DBLP:journals/corr/abs-2009-10297} and \emph{CrystalBLEU}~\cite{DBLP:conf/kbse/EghbaliP22} attempt to address these limitations by introducing additional heuristics such as AST matching.

The metric \emph{pass@k} reports the likelihood that a model correctly completes a code snippet at least once within $k$ attempts. Originally introduced by \citeauthor{DBLP:conf/nips/KulalPC0PAL19}~\cite{DBLP:conf/nips/KulalPC0PAL19} as the success rate within a budget of $B$ program attempts (the role now played by $k$ in \emph{pass@k}) and popularized by \citeauthor{DBLP:journals/corr/abs-2107-03374}~\cite{DBLP:journals/corr/abs-2107-03374}, it ranges from 0 (no correct solution in $k$ tries) to 1 (at least one correct solution), with correctness typically defined by passing test cases.

The \emph{pass@k} metric is defined as:

\[
\text{pass@}k = 1 - \frac{\binom{n-c}{k}}{\binom{n}{k}}\,, \text{where:}
\]

\begin{itemize}
  \item $n$ is the total number of generated samples per prompt,
  \item $c$ is the number of correct samples among $n$, and
  \item $k$ is the number of samples drawn.
\end{itemize}

The choice of $k$ depends on the downstream task: \emph{pass@1} is critical for single-suggestion scenarios such as code completion, while higher $k$ values (e.g., 2, 5, 10) assess multi-attempt capability and are commonly reported in technical reports of recent code LLMs (e.g., Code Llama~\cite{DBLP:journals/corr/abs-2308-12950}, DeepSeek-Coder~\cite{DBLP:journals/corr/abs-2401-14196}, Qwen2.5-Coder~\cite{DBLP:journals/corr/abs-2409-12186}, StarCoder~\cite{DBLP:journals/tmlr/LiAZMKMMALCLZZW23}).
However, \emph{pass@k} is not a universal metric suitable for all generation tasks.
It requires a binary notion of correctness, making the metric appropriate for code synthesis evaluated via unit tests, but unsuitable for open-ended generation tasks such as comment generation, where multiple valid outputs exist.

A complementary perspective from industry practice is \emph{pass\textsuperscript{k}} (also written \emph{pass\^{}k}), which measures the probability that \emph{all} $k$ trials succeed rather than \emph{at least one}~\cite{Anthropic2025evals}. While \emph{pass@k} increases with $k$, \emph{pass\textsuperscript{k}} decreases, making it useful for assessing reliability in deployment scenarios where consistent success matters (e.g., customer-facing agents).

If a study evaluates an LLM-based tool for supporting humans, a relevant metric is the acceptance rate, meaning the ratio of all accepted artifacts (e.g., test cases, code snippets) in relation to all artifacts that were generated and presented to the user.
Another way of evaluating LLM-based tools is calculating inter-model agreement, which reveals how much a tool's performance depends on specific models and versions.
Metrics used to measure inter-model agreements include general agreement (percentage), \emph{Cohen's} $\kappa$, and \emph{Krippendorff's} $\alpha$ (see \humanvalidation for recommended thresholds and best practices for measuring agreement).

Common problem types for LLM-based studies are classification, recommendation, and generation, each requiring different metrics~\cite{DBLP:journals/tosem/HouZLYWLLLGW24}. \citeauthor{DBLP:journals/corr/abs-2505-08903}~\cite{DBLP:journals/corr/abs-2505-08903} categorized 191 LLM benchmarks by SE task, providing a valuable reference. For an overview of code foundation models, agents, and their evaluation, see~\citeauthor{yang2025code}~\cite{yang2025code}. From a practitioner perspective, \citeauthor{Anthropic2025evals}~\cite{Anthropic2025evals} categorize evaluation approaches for LLM-based agents into code-based graders (e.g., test suite execution, static analysis), model-based graders (e.g., rubric-scored LLM judgments), and human graders. This taxonomy may help researchers systematically design evaluation strategies for agent-based tools. Common metrics include \emph{BLEU}, \emph{pass@k}, \emph{Accuracy@k}, and \emph{Exact Match} for generation; \emph{Mean Reciprocal Rank} for recommendation; and \emph{Precision}, \emph{Recall}, \emph{F1-score}, and \emph{Accuracy} for classification.

\paragraph{Benchmark Examples:}

Benchmarks used for code generation include \emph{HumanEval} (available on \href{https://github.com/openai/human-eval}{GitHub}) \cite{DBLP:journals/corr/abs-2107-03374}, \emph{MBPP} (available on \href{https://huggingface.co/datasets/google-research-datasets/mbpp}{Hugging Face}) \cite{DBLP:journals/corr/abs-2108-07732},
\emph{ClassEval} (available on \href{https://github.com/FudanSELab/ClassEval}{GitHub}) \cite{DBLP:conf/icse/Du0WWL0FS0L24}, \emph{LiveCodeBench} (available on \href{https://github.com/LiveCodeBench/LiveCodeBench}{GitHub}) \cite{DBLP:journals/corr/abs-2403-07974}, and \emph{SWE-bench} (available on \href{https://github.com/swe-bench/SWE-bench}{GitHub}) \cite{DBLP:conf/iclr/JimenezYWYPPN24}.
An example of a code translation benchmark is \emph{TransCoder}~\cite{DBLP:conf/nips/RoziereLCL20} (available on \href{https://github.com/facebookresearch/CodeGen}{GitHub}).
\citeauthor{DBLP:journals/corr/abs-2601-11895}'s \emph{DevBench} (available on \href{https://github.com/microsoft/devbench}{GitHub})~\cite{DBLP:journals/corr/abs-2601-11895} applies the synthesis strategy, with 1,800 code completion instances derived from developer telemetry.
% GROUNDING: 1,800 instances; ``categories \dots\ derived from analysis of over one billion real developer interactions'' (Section 2.1).

Most benchmarks focus on generation tasks, but benchmarks for classification and recommendation also exist.
For classification, \emph{DiverseVul} (available on \href{https://github.com/wagner-group/diversevul}{GitHub})~\cite{DBLP:conf/raid/0001DACW23} provides vulnerable and non-vulnerable functions evaluated using standard classification metrics.
For recommendation, \emph{CodeSearchNet} (available on \href{https://github.com/github/CodeSearchNet}{GitHub})~\cite{DBLP:journals/corr/abs-1909-09436} contains code-documentation pairs evaluated using \emph{Mean Reciprocal Rank}.

\guidelinesubsubsection{Benefits}

Benchmarks and metrics improve reproducibility and comparability across studies, leading to faster improvements in the cumulative body of knowledge. It is simply easier to assess a new system against one or a few benchmarks than hundreds of competing systems, and when most systems are assessed against the same benchmarks and metrics, we can see which approaches work best (at least, on those benchmarks). This comparison enables progress tracking; for example, when researchers iteratively improve a new LLM-based tool and test it against benchmarks after substantial changes. For practitioners, leaderboards (published benchmark results of models) support selecting models for downstream tasks. Meanwhile, baselines help us understand just how much LLMs improve performance over non-LLM-enabled alternatives. 

\guidelinesubsubsection{Challenges}

Computer scientists commonly use simple metrics (e.g., lines of code, CPU time) as proxies for complex, multidimensional latent variables (e.g., system size, environmental sustainability), without empirically validating that the metrics capture the intended construct~\cite{DBLP:conf/ease/RalphT18}. Unlike fields such as psychology where measurement theory has a longer tradition~\cite{borsboom2005measuring}, many AI metrics and benchmarks lack both theoretical underpinnings and empirical validation of their construct, measurement, and ecological validity.

These validity issues are widespread, not isolated. In a survey of 572 code benchmarks released between 2014 and 2025, \citeauthor{cao2025should}~\cite{cao2025should} found that 84.2\% did not consider test-suite coverage when constructing test cases, 64.0\% reported single-pass evaluations without controlling for randomness, and 82.5\% did not address data contamination. \citeauthor{bean2025measuring}~\cite{bean2025measuring} reported comparable patterns in a parallel review of 445 LLM benchmarks from ML and NLP venues. As a result, model performance on these benchmarks can reflect benchmark artifacts rather than the capability the benchmark claims to test.

For example, \emph{pass@k} is intended to measure a model's \textit{functional correctness}, that is, its capability to generate code that produces the expected output for a given specification. However, functional correctness is only one dimension of code quality. \emph{pass@k} does not capture maintainability, readability, security, or efficiency of the generated code, all of which are critical for downstream use. Furthermore, ``correctness'' is defined entirely by test suites, whose coverage is itself unvalidated: a solution that passes all provided tests may still be incorrect for untested inputs. Whether a test suite adequately operationalizes correctness for a given task is a subjective judgment that is rarely examined.

Similarly, \emph{HumanEval} is intended to measure a model's capability to \textit{synthesize short Python functions from docstring specifications}. This operationalizes a narrow slice of ``code generation capability'': it covers neither multi-file tasks, nor debugging, nor the use of existing codebases. These tasks dominate real-world software engineering~\cite{Chandra2025benchmarks, yang2025code}. Notably, neither \emph{pass@k} nor \emph{HumanEval} has undergone rigorous empirical validation of its construct validity; their widespread adoption rests on face validity and convenience rather than on evidence that they reliably measure what researchers intend them to measure~\cite{cao2025should}. \emph{HumanEval} also contains implementation, documentation, and test-case bugs in its original release, which directly affect score interpretation~\cite{DBLP:conf/nips/LiuXW023}; comparable issues have been documented in other widely-used code benchmarks~\cite{cao2025should}.

Benchmarks and metrics also have generalizability problems. Prominent LLM benchmarks such as \emph{HumanEval} and \emph{MBPP} use Python, so researchers can optimize for Python's idiosyncrasies, producing gains that do not generalize to other languages. Similarly, the metric \emph{BLEU-N} is a syntactic metric, so code can score highly without being executable. The metric \emph{Exact Match}, meanwhile, does not account for functional equivalence of syntactically different code. Both \emph{BLEU-N} and \emph{Exact Match} are influenced by code formatting, which confounds their intended use.

Execution-based metrics such as \emph{pass@k} directly evaluate correctness by running test cases, but they require an execution environment.
When metric values look surprising, examine the specific test cases driving them: common causes include outliers, bugs in test suites or scoring code, and items that exercise capabilities other than the one being measured.

Finally, benchmark data contamination, where the benchmark itself is part of the training data, may lead to artificially high performance if the model remembers the solution from the training data rather than deriving it from the input~\cite{DBLP:journals/corr/abs-2406-04244}. Therefore, for proprietary LLMs that do not release their training data, researchers \should consider using human validation, curating new data, or refactoring existing data~\cite{DBLP:journals/corr/abs-2406-04244}. 

To mitigate contamination, researchers can create new benchmark datasets by collecting data after a specified cutoff date; researchers \must disclose data sources and collection dates for each release. However, this temporal approach requires continuous updates as new models may include the benchmark in future training data. Alternatively, keeping the benchmark private prevents inclusion in training sets, but requires trust in the benchmark creator and a system to execute it without leaking data. A third strategy is to synthesize benchmark instances rather than draw them from public sources, which avoids training-data contamination at the cost of requiring an instance generator and a validation pipeline. Researchers can also evaluate how contaminated their benchmark is~\cite{DBLP:journals/corr/abs-2502-00678}.


\guidelinesubsubsection{Study Types}

This guideline \must be followed for all study types that automatically evaluate the performance of LLMs or LLM-based tools.
The design of a benchmark and the selection of appropriate metrics are highly dependent on the specific study type and research goal.
Recommending specific metrics for specific study types is beyond the scope of these guidelines, but \citeauthor{DBLP:journals/corr/abs-2505-08903} provide a good overview of existing metrics for evaluating LLMs~\cite{DBLP:journals/corr/abs-2505-08903}.

For \benchmarkingtasks, this guideline is of primary importance: researchers \must use established benchmarks or rigorously justify the creation of new ones, \must report standard metrics to enable cross-study comparison, and \should report an error analysis alongside the primary scores so that readers can assess whether reported gains reflect the target capability or other capabilities the task happens to require.
For \annotators, the research goal might be to assess which model comes close to a ground truth dataset created by human annotators.
Especially for open annotation tasks, selecting suitable metrics to compare LLM-generated and human-generated labels is important.
In general, annotation tasks can vary widely.
Are multiple labels allowed for the same sequence? Are the available labels predefined, or should the LLM generate a set of labels independently?
Due to this task dependence, researchers \must justify their metric choice, explaining what aspects of the task it captures together with known limitations.
For \newtools, researchers \should benchmark the tool against suitable baselines using appropriate metrics that capture the tool's intended contribution.
For \judges, researchers \should report inter-rater agreement metrics and validity measures to demonstrate the reliability and quality of LLM judgments.
For \synthesis, researchers \should specify metrics for comparing synthesized outputs (e.g., coverage, faithfulness) and justify their appropriateness for the synthesis task.
For \llmusage, researchers \should justify the measurement instruments and metrics used for studying LLM usage patterns, including any survey scales or behavioral measures.
For \subjects, researchers \should compare simulated and real human responses using appropriate metrics to assess simulation fidelity.
If researchers assess a well-established task such as code generation, they \should report standard metrics such as \emph{pass@k} and compare the performance between models.
If non-standard metrics are used, researchers \must state their reasoning.

\guidelinesubsubsection{Advice for Reviewers}

Reviewers should expect manuscripts to:
(1) clearly identify the constructs or variables the study aims to measure (e.g., LLM performance, quality of generated code), including independent, dependent, and control variables;
(2) present their measurement model, i.e., which metrics, benchmarks, or baselines are used and how they relate to the target constructs;
(3) justify the selection of metrics, benchmarks, and baselines;
(4) discuss \textit{in detail} the assumptions, reliability, and validity (\textit{especially construct validity}) of each benchmark and metric;
(5) articulate any limitations regarding construct and measurement validity.

As with other guidelines, missing information about baselines or metrics is typically a revision request. However, vague descriptions that conflate broad concerns (e.g., effectiveness, quality) with specific counting methods should be questioned. Ubiquity of a benchmark or metric does not imply validity or appropriateness for a given context. Manuscripts should convey a solid understanding of construct and measurement validity by explaining and justifying their measurement models.

\guidelinesubsubsection{See Also}
\begin{itemize}[label=$\rightarrow$]
    \item \seeopenllm: Using an open LLM as a baseline supports reproducible benchmark comparisons across studies.
    \item \seehumanvalidation: Human evaluation captures aspects of LLM outputs that automated metrics cannot measure.
    \item \seelimitationsmitigations: Every benchmark and metric choice carries construct-validity threats that authors must discuss.
\end{itemize}